\appendixchapter{А}{Техническое задание}

\begin{center}
\textbf{ТЕХНИЧЕСКОЕ ЗАДАНИЕ}\\
на разработку информационной системы\\
для создания и управления туристическими маршрутами\\
«Guide Helper»\\
\vspace{0.5em}
ЕСПД. ГОСТ 19.201-78~\cite{gost19201}
\end{center}

\section{Введение}

Настоящее техническое задание распространяется на разработку информационной системы для создания и управления туристическими маршрутами с геопривязанными фотографиями «Guide Helper».

Условное обозначение системы: Guide Helper.

Плановые сроки начала и окончания работ: сентябрь 2025~г.~--- апрель 2026~г.

\section{Основания для разработки}

Основанием для разработки является задание на выполнение выпускной квалификационной работы студента кафедры программной инженерии.

Наименование работы: Разработка информационной системы для создания туристических маршрутов с геопривязанными фотографиями.

\section{Назначение разработки}

Функциональное назначение: автоматизация деятельности туристических гидов и инструкторов по созданию, визуализации и распространению туристических маршрутов с мультимедийным контентом.

Эксплуатационное назначение: система предназначена для использования туристическими гидами, инструкторами, организаторами походов и туристами, планирующими самостоятельные маршруты.

\section{Требования к программе}

\subsection{Требования к функциональным характеристикам}

Система должна обеспечивать следующие функциональные возможности:

\begin{enumerate}
  \item Регистрация, авторизация и управление профилем пользователей с ролевой моделью (пользователь, модератор, администратор).
  \item Создание туристических маршрутов с точками на интерактивной карте; прикрепление геопривязанных фотографий к точкам маршрута.
  \item Автоматическое построение маршрута между точками через сервис маршрутизации OSRM.
  \item Расчёт статистик маршрута: расстояние (формула Гаверсинуса), набор высоты, расчётное время прохождения (формула Нейсмита), автоматическая классификация сложности.
  \item Асинхронная серверная обработка загруженных фотографий: сжатие, масштабирование, генерация миниатюр; хранение в объектном хранилище.
  \item Публикация маршрутов в общедоступном каталоге; поиск, фильтрация и сортировка маршрутов.
  \item Публичный просмотр опубликованных маршрутов по уникальной ссылке без регистрации.
  \item Социальные функции: комментарии, отметки «нравится», рейтинги (1--5~звёзд), закладки.
  \item ИИ-ассистент с поддержкой нескольких провайдеров языковых моделей и вызовом инструментов (геокодирование, поиск маршрутов, прогноз погоды).
  \item Уведомления о социальных событиях через WebSocket в реальном времени.
  \item Импорт маршрутов из форматов GeoJSON; экспорт в форматы GPX и KML.
  \item Административная панель: управление пользователями, модерация контента, управление категориями, просмотр статистики.
  \item Прогрессивное веб-приложение и десктоп-приложение на основе единой кодовой базы.
\end{enumerate}

\subsection{Требования к надёжности}

\begin{itemize}
  \item Доступность системы: не менее 99\% времени в год.
  \item Время отклика API: не более 500~мс для 95-го перцентиля запросов.
  \item Время обработки фотографий: не более 10~секунд для файлов размером до 10~МБ.
  \item Устойчивость к отказам: сбой одного микросервиса не должен приводить к полной недоступности системы.
  \item Гарантированная доставка задач обработки фотографий: повторная обработка при сбое потребителя.
\end{itemize}

\subsection{Требования к условиям эксплуатации}

\begin{itemize}
  \item Температура окружающей среды для серверного оборудования: от +10 до +35~°C.
  \item Относительная влажность: 20--80\% без конденсата.
  \item Непрерывный круглосуточный режим эксплуатации (24/7).
\end{itemize}

\subsection{Требования к составу и параметрам технических средств}

Минимальные требования к серверу:
\begin{itemize}
  \item процессор: 4~ядра, тактовая частота 2~ГГц и выше;
  \item оперативная память: 8~ГБ и более;
  \item дисковое пространство: 50~ГБ и более (SSD);
  \item сетевое подключение: 100~Мбит/с.
\end{itemize}

Минимальные требования к клиентскому рабочему месту:
\begin{itemize}
  \item современный веб-браузер с поддержкой ES2020 (Chrome~90+, Firefox~89+, Safari~14+, Edge~90+);
  \item разрешение экрана: 1280$\times$720 пикселей и выше;
  \item сетевое подключение: 5~Мбит/с.
\end{itemize}

\subsection{Требования к информационной и программной совместимости}

\begin{itemize}
  \item Серверная часть развёртывается как набор Docker-контейнеров; поддерживаются Linux (x86\_64, arm64), платформа Docker Engine 24+.
  \item Взаимодействие клиента и сервера: REST API через HTTP/HTTPS, WebSocket, Server-Sent Events.
  \item Формат обмена данными: JSON; идентификаторы ресурсов: UUID v4.
  \item Интеграция с внешними сервисами: OpenStreetMap / Nominatim, Open-Meteo Elevation API, OSRM, AI-провайдеры и OpenAI-совместимые API (опционально).
  \item Форматы импорта/экспорта маршрутов: GeoJSON, GPX, KML.
\end{itemize}

\section{Требования к программной документации}

В состав программной документации должны входить:
\begin{itemize}
  \item техническое задание (настоящий документ);
  \item пояснительная записка к выпускной квалификационной работе;
  \item руководство пользователя (часть пояснительной записки);
  \item программа и методика испытаний (Приложение Б).
\end{itemize}

\section{Технико-экономические показатели}

Ориентировочная трудоёмкость разработки: 570~человеко-часов (11~этапов, подробный расчёт в разделе~3.5 пояснительной записки).

Расчётная стоимость разработки: 1~897~508~рублей (включая фонд оплаты труда, страховые взносы, накладные расходы и амортизацию).

Ожидаемый экономический эффект: сокращение времени на создание и публикацию одного маршрута с~2,5~часа до~30~минут; срок окупаемости~--- 0,8~года при целевой аудитории 2~000~гидов.

Программа разрабатывается как проприетарное программное обеспечение. Права на
исходный код сохраняются за разработчиком, а использование системы предполагается
в виде централизованно сопровождаемого сервиса.

\section{Стадии и этапы разработки}

\begin{spacing}{1.0}
\small
\begin{longtable}{|c|p{8cm}|c|}
\caption{Стадии и этапы разработки}
\label{tab:stages} \\
\hline
\textbf{№} & \textbf{Наименование этапа} & \textbf{Срок} \\
\hline
\endfirsthead
\multicolumn{3}{r}{\small Продолжение таблицы~\ref{tab:stages}} \\
\hline
\endhead
\longtablecontinued{3}{tab:stages}
\endfoot
\longtableended{3}{tab:stages}
\endlastfoot
1 & Анализ предметной области и конкурентов & Сентябрь 2025 \\
\hline
2 & Формирование требований, проектирование архитектуры & Октябрь 2025 \\
\hline
3 & Проектирование базы данных и интерфейса & Октябрь 2025 \\
\hline
4 & Реализация сервиса аутентификации & Ноябрь 2025 \\
\hline
5 & Реализация сервиса маршрутов & Ноябрь--декабрь 2025 \\
\hline
6 & Реализация асинхронной обработки фотографий & Декабрь 2025 \\
\hline
7 & Реализация клиентского веб-приложения & Январь 2026 \\
\hline
8 & Интеграция ИИ-ассистента & Февраль 2026 \\
\hline
9 & Настройка автоматической сборки, проверки и Kubernetes-развёртывания & Февраль--март 2026 \\
\hline
10 & Тестирование и устранение дефектов & Март 2026 \\
\hline
11 & Оформление документации & Март--апрель 2026 \\
\hline
\end{longtable}
\end{spacing}

\section{Порядок контроля и приёмки}

Приёмка системы осуществляется в соответствии с Программой и методикой испытаний (Приложение~Б). Критерии приёмки:
\begin{itemize}
  \item успешное прохождение всех тест-кейсов функционального тестирования;
  \item выполнение нефункциональных требований по времени отклика и доступности;
  \item отсутствие критических и блокирующих дефектов;
  \item наличие полного комплекта программной документации.
\end{itemize}
